% Abstract für Bachelorarbeit
% Nach Schreibwerkstatt Universität Stuttgart: Prägnant, objektiv, ohne Ausblick, ohne Vollformen

\chapter*{Abstrakt}
\markboth{Abstrakt}{Abstrakt}
\thispagestyle{scrheadings}
\addcontentsline{toc}{chapter}{Abstrakt}

Klassische \gls{gls-SIP}-Telefonsysteme ermöglichen Echtzeitkommunikation, jedoch bieten sie keine intelligente Sprachverarbeitung. Diese Arbeit untersucht die technische Machbarkeit der Integration von Large Language Models in containerbasierte \gls{gls-SIP}-Telefonsysteme. Die zentrale Forschungsfrage lautet: Unter welchen Bedingungen ermöglichen containerbasierte Ansätze eine stabile Verbindung von \gls{gls-SIP}-basierter Telefonie mit \gls{gls-KI}-Sprachmodellen?

Der Prototyp verbindet drei Kernkomponenten: Asterisk übernimmt die \gls{gls-SIP}-\\Signalisierung, LiveKit verarbeitet \gls{gls-WebRTC}-Medienströme, und ein Python-Agent steuert die \gls{gls-KI}-Logik. Die Sprachverarbeitung erfolgt über OpenAI Whisper für die Spracherkennung, GPT-4o-mini für die Dialogverwaltung und OpenAI Text-to-Speech für die Sprachsynthese. Die Docker-basierte Implementierung garantiert Reproduzierbarkeit. Die iterative Entwicklung begann mit einem leichtgewichtigen \gls{gls-SIP}-Agenten für erste Schritte und das Erlernen grundlegender \gls{gls-SIP}-Mechanismen und führte schrittweise zur finalen Architektur.

Die Evaluation konzentriert sich auf die \gls{gls-SIP}-Integration. Fünf repräsentative Testanrufe wurden detailliert ausgewertet. Die Median-Latenz beträgt 0,78 s für Spracherkennung und 1,41 s für Sprachsynthese. Das System zeigt stabile bidirektionale Audioübertragung ohne Timeouts. Eine \gls{gls-STT}-Fehlinterpretation wurde dokumentiert. Die CPU-Auslastung liegt bei 11,2\%, der RAM-Bedarf zwischen 11,2 und 12,1 GB. Die Audioqualität erreicht \gls{gls-MOS}-Werte von 4,0 bis 4,5.

Die Arbeit präsentiert eine \gls{gls-ETSI}-konforme Architektur, demonstriert eine funktionsfähige Implementierung, liefert Evaluationsdaten aus \gls{gls-SIP}-Tests und formuliert Handlungsempfehlungen für kleine und mittlere Unternehmen. Open-Source-Lösungen ermöglichen Datenschutz-konforme Deployments ohne Vendor-Lock-in. Die Ergebnisse zeigen: \gls{gls-KI}-Sprachmodelle lassen sich ohne spezialisierte Hardware mit klassischen \gls{gls-VoIP}-Systemen kombinieren, um dialogorientierte Sprachverarbeitung zu ermöglichen.\\

\noindent\textbf{Schlagwörter:} Large Language Model, \gls{gls-SIP}, Telefonie, Docker-basierte Architektur, Asterisk, LiveKit, \gls{gls-KI}-Integration
